\documentclass[xcolor=dvipsnames,t,aspectratio=169]{beamer}

\usecolortheme{rose}
\usecolortheme{dolphin}
\usetheme{Boadilla}

\input{../../templates/slides/imports}
\input{../../templates/slides/settings}
\input{../../templates/slides/commands}

\titlegraphic{
    \includegraphics[scale = 0.5]{../../templates/slides/logo}
}

\logo{
\begin{tikzpicture}[overlay,remember picture]
\node[xshift=-.75cm, yshift=-1cm] at (current page.30){
    \includegraphics[width=0.1\textwidth]{../../templates/slides/logo}};
\end{tikzpicture}
}
\usetikzlibrary{positioning}

\newcommand{\highlight}[1]{{\color{nes_dark_orange} #1}}

\title{NumPy e computação numérica }

\author{Eduardo Adame}

\date{{\color{nes_dark_purple}\textbf{Ciência de Dados}\\[0.5em] 20 de março de 2026 }}

\begin{document}

\frame[plain]{\titlepage}
\setcounter{framenumber}{0}


\section{NumPy}

\begin{frame}[c, fragile]{Por que NumPy?}

Problema: Python puro não é eficiente para computação numérica.

\begin{code}[Python vanilla]{python}
L = list(range(1000))
[i**2 for i in L]
\end{code}

\begin{code}[NumPy]{python}
import numpy as np
a = np.arange(1000)
a**2
\end{code}

\begin{itemize}
\item NumPy é \highlight{vetorizado}
\item Executa operações em \highlight{C por baixo}
\item Muito mais rápido e eficiente
\end{itemize}

\end{frame}


\begin{frame}[c, fragile]{O que é um array?}
\begin{itemize}
\item Estrutura de dados homogênea
\item Elementos do mesmo tipo
\item Armazenamento contíguo em memória
\end{itemize}

\begin{code}[Exemplo]{python}
a = np.array([1,2,3])
a.shape
a.dtype
\end{code}

\begin{display}[Importante]
NumPy não é uma lista melhor — é outra estrutura.
\end{display}

\end{frame}

\begin{frame}[c, fragile]{Criando arrays}

\begin{code}[Criação]{python}
np.array([1,2,3])

np.arange(0,10)

np.linspace(0,1,5)

np.zeros((3,3))

np.ones((2,2))

np.eye(3)
\end{code}

\begin{itemize}
\item Evite construir arrays manualmente
\item Prefira funções vetorizadas
\end{itemize}

\end{frame}

\begin{frame}[c, fragile]{Shape e dimensão}

\begin{code}[Exemplo]{python}
a = np.array([[1,2,3],[4,5,6]])

a.shape   # (2,3)
a.ndim    # 2
\end{code}

\begin{itemize}
\item shape = tamanho em cada dimensão
\item ndim = número de dimensões
\end{itemize}

\end{frame}

\begin{frame}[c, fragile]{Indexação}

\begin{code}[Indexação]{python}
a = np.arange(10)

a[0]
a[-1]
a[2:5]
a[::2]
\end{code}

\begin{code}[2D]{python}
b = np.array([[1,2],[3,4]])

b[0,1]
\end{code}

\end{frame}

\begin{frame}[c, fragile]{Operações vetorizadas}

\begin{code}[Elementwise]{python}
a = np.array([1,2,3])

a + 1
a * 2
a ** 2
\end{code}

\begin{itemize}
\item Sem loops explícitos
\item Muito mais rápido
\end{itemize}

\end{frame}

\begin{frame}[c, fragile]{Comparação com Python}

\begin{code}[Python]{python}
[i+1 for i in range(10000)]
\end{code}

\begin{code}[NumPy]{python}
a = np.arange(10000)
a + 1
\end{code}


\highlight{Evite} loops sempre que possível.


\end{frame}

\begin{frame}[c, fragile]{Broadcasting}

\begin{code}[Exemplo]{python}
a = np.array([[1,2,3],
              [4,5,6]])

b = np.array([1,2,3])

a + b
\end{code}

\begin{itemize}
\item NumPy expande dimensões automaticamente
\item Evita loops explícitos
\end{itemize}

\end{frame}

\begin{frame}[c, fragile]{Reduções}

\begin{code}[Exemplo]{python}
a = np.array([[1,2],[3,4]])

a.sum()
a.sum(axis=0)
a.mean()
a.std()
\end{code}

\begin{itemize}
\item axis controla dimensão da operação
\end{itemize}

\end{frame}

\begin{frame}[c, fragile]{Operações booleanas e filtros}

\begin{code}[Comparações]{python}
a = np.array([1,2,3,4])

a > 2
a == 3
\end{code}

\begin{code}[Máscaras booleanas]{python}
mask = a > 2
a[mask]
\end{code}

\begin{code}[Filtro direto]{python}
a[a > 2]
\end{code}

\end{frame}

\begin{frame}[c, fragile]{Operações booleanas e filtros}

\begin{code}[Combinações]{python}
(a > 1) & (a < 4)
(a == 1) | (a == 4)
\end{code}

\begin{itemize}
\item Permite selecionar subconjuntos de dados
\item Base para \highlight{data filtering} em ciência de dados
\item Evita loops explícitos
\end{itemize}

\end{frame}

\begin{frame}[c, fragile]{Álgebra linear}

\begin{code}[Produto]{python}
A @ B
\end{code}

\begin{itemize}
\item \texttt{*} não é produto matricial
\item use \texttt{@} ou \texttt{np.dot}
\end{itemize}

\end{frame}

\begin{frame}[c, fragile]{Performance}

\begin{code}[Benchmark]{python}
import time

initial_time = time.time()
# python
python_time = time.time() - initial_time
initial_time = time.time()
# numpy
numpy_time = time.time() - initial_time
\end{code}

\begin{itemize}
\item NumPy usa memória contígua
\item operações vetorizadas
\item evita overhead do Python
\end{itemize}

\end{frame}


\begin{frame}[c, noframenumbering, plain]
    \frametitle{~}
        \vfill
        \begin{center}
            {\Huge Obrigado!}\vspace{1.5em}\\
            {\Large \highlight{Alguma dúvida?}}\\
        \end{center}
        \vfill
        \begin{center}
            {\small Agora vamos para os exercícios!}
        \end{center}
\end{frame}

\end{document}